\chapter{CrewAI as a Writing Organization}
CrewAI expresses a useful production pattern: complex intellectual work can be divided among agents with specialized roles. In this project the crew contains a planner, researcher, writer, editor, LaTeX engineer, and QA reviewer. The purpose is not to pretend that agents are people. The purpose is to force the workflow to expose responsibilities, expected outputs, and quality gates.

\section{Roles and Tasks}
\begin{table}[h]
\centering
\begin{tabular}{lll}
\toprule
Role & Main responsibility & Output \\
\midrule
Planner & Decompose assignment & Outline and checklist \\
Researcher & Gather claims and references & Evidence notes \\
Writer & Draft chapters & Markdown manuscript \\
Editor & Improve clarity and caution & Reviewed manuscript \\
LaTeX Engineer & Typeset document & TeX files and figures \\
QA Reviewer & Validate requirements & Final checklist \\
\bottomrule
\end{tabular}
\caption{CrewAI-style agent team for the assignment.}
\end{table}

The most important design choice is the context chain. The writer should not invent from nothing; it should consume the research output. The editor should not rewrite the assignment; it should improve the writer's artifact. The LaTeX engineer should not decide the argument; it should typeset an approved manuscript.
